\section{Institutions and the Information Required for Procurement}
\label{sec:r19-information}
A value oracle is useful only relative to the decision for which it is queried. In particular, identifying each response coordinate separately need not identify the return to buying their joint production. The following result makes this distinction exact. It concerns a purchaser who observes values, not a purchaser who has already been given the transition and reward arrays. Known dynamics provide additional information and are treated separately below.

Let $k\geq2$. A dynamic contract has a compact set of attainable response vectors $Z\subset[0,1]^k$, including public randomizations, and private value
\begin{equation}
 W_Z(t)=\max_{z\in Z}t\cdot z,\qquad t\in\mathbb R^k.
 \label{eq:r19-support-oracle}
\end{equation}
At the unperturbed contract all responses are privately optimal. The purchaser receives $w\cdot z$, where $w\in[0,\infty)^k$ describes the institution's allocation of service, receipts, or quality. A safe offer yields $c$ independently of response. The return to the response-sensitive offer is its worst equilibrium response, $r(Z,w)=\min_{z\in Z}w\cdot z$. The relevant information problem is to select between these offers, rather than to select the agent's preferred response. Perturbations $t$ are queries used to evaluate this fixed contract; they are not imposed on its eventual recipient.

\begin{theorem}[Institution-dependent minimal information]
\label{thm:r19-information}
Set $Z_1=\operatorname{co}\{e_1,\ldots,e_k\}$ and $Z_0=\operatorname{co}(Z_1\cup\{0\})$, and let $\rho=\min_i w_i>0$ and $0<c<\rho$.
\begin{enumerate}
\item Every query in $\mathcal T=\{t:\max_i t_i\geq0\}$ has the same exact value in the two economies. Consequently, even an unlimited sequence of coordinate queries, with directions and magnitudes chosen adaptively, cannot identify the optimal procurement decision. The response-sensitive offer is strictly preferred in economy 1 and the safe offer in economy 0.
\item The minimax regret of a randomized decision based solely on these messages is
\begin{equation}
 \mathcal R^*=\frac{c(\rho-c)}{\rho}>0.
 \label{eq:r19-regret}
\end{equation}
This is an information loss, not a numerical approximation error.
\item A single additional query $t=-hw$, $h>0$, identifies the optimal offer under exact values. If the query is reported with an adversarial absolute error at most $\varepsilon$, identification for every permissible report is possible if and only if $h\rho>2\varepsilon$. When $h\rho\leq2\varepsilon$, the minimax regret remains \eqref{eq:r19-regret}.
\item With prior probability $p$ on economy 1, the gross expected value of the exact joint query is
\begin{equation}
 p\rho+(1-p)c-\max\{p\rho,c\}.
 \label{eq:r19-information-value}
\end{equation}
If any component of $w$ is zero, both economies instead have robust return zero and this distinction has no value for a positive-return safe offer.
\end{enumerate}
Both economies can be implemented by a finite-horizon dynamic contract with a first-date choice followed by an absorbing operating state. The result therefore holds within the class of dynamic response problems, not merely for an abstract oracle unrelated to a policy problem.
\end{theorem}

\begin{proof}
Linearity gives $W_{Z_1}(t)=\max_i t_i$ and $W_{Z_0}(t)=\max\{0,\max_i t_i\}$. These expressions coincide throughout $\mathcal T$. Every coordinate vector $se_i$ belongs to $\mathcal T$, including when $s<0$, because $k\geq2$. Under an adaptive query rule, identical past messages lead to the same distribution of subsequent queries and messages; induction gives identical transcripts. On the other hand, $r(Z_1,w)=\rho$ and $r(Z_0,w)=0$.

Write $a\in[0,1]$ for the probability of selecting the response-sensitive offer after a common transcript. Its regrets in economies 1 and 0 are respectively $(1-a)(\rho-c)$ and $ac$. Their maximum is minimized at $a=(\rho-c)/\rho$, where both equal \eqref{eq:r19-regret}. Randomizing query choices cannot improve this value because their transcripts have the same law in both economies.

The joint query has values $-h\rho$ and zero. The two sets of permissible reports are the closed intervals $[-h\rho-\varepsilon,-h\rho+\varepsilon]$ and $[-\varepsilon,\varepsilon]$. They are disjoint exactly when $h\rho>2\varepsilon$. If they intersect, an adversarial oracle can give the same report in either economy. The previous regret lower bound follows, and ignoring that report attains the bound. Under exact information the purchaser obtains $\rho$ in economy 1 and $c$ in economy 0; without the information its best expected return is $\max\{p\rho,c\}$. This proves \eqref{eq:r19-information-value}. A zero weight gives a zero-return vertex in both economies. Finally, assign response $e_i$ to each initial action, give every action zero private payoff at the center, and add the zero-response action only in economy 0. An initial public lottery generates precisely the stated convex hulls. Further dates with zero additional payoff and response embed the construction in any longer finite horizon.
\end{proof}

The theorem identifies an exact boundary on useful information: the complete value function on $\mathcal T$ is insufficient, whereas one purchaser-directed query outside that set suffices. The claim is not that a joint query solves every dynamic contract. It is that coordinate precision cannot substitute for the missing joint direction even in a dynamically implementable problem, and that the institution determines both the direction and the value of acquiring it. Fee incidence is one example: when a purchaser begins to receive an additional output, an initially irrelevant response coordinate can become necessary for the ranking. The theorem concerns the stated information experiment; it does not claim a universal lower bound against an algorithm that reads the model's transition arrays.

\subsection{Exact messages and general near-optimal responses}
The preceding construction can be related directly to the paper's finite-query response set. Suppose $W(t_i)=v_i$ is known exactly at finitely many directions, including $t_0=0$, and the observations are consistent with at least one affine-response economy whose queried optima are attained. Features lie in a known compact polytope $K$. Define
\begin{equation}
 C_\eta=\{z\in K:t_i\cdot z\leq v_i-v_0+\eta\text{ for every }i\}.
 \label{eq:r19-exact-messages}
\end{equation}

\begin{proposition}[Sharp exact-message response set]
\label{prop:r19-sharp}
Across affine-response economies consistent with the exact messages, $C_\eta$ is exactly the set of features that can occur as an attained $\eta$-optimal response at the center. Thus a strictly positive lower bound on $w\cdot z-c$ over $C_\eta$ identifies robust superiority of the response-sensitive offer. If the minimum is nonpositive, such strict superiority is not implied by the messages alone.
\end{proposition}
\begin{proof}
If a response has center payoff $a\geq v_0-\eta$, then $a+t_i\cdot z\leq v_i$, giving \eqref{eq:r19-exact-messages}. Conversely, take any consistent economy and any $z\in C_\eta$. Add one policy with payoff $v_0-\eta+t\cdot z$. At every query its payoff is no greater than $v_i$, and at the center it is no greater than $v_0$. Hence all messages remain unchanged and the added policy is $\eta$-optimal. Randomizing with existing policies does not raise their maximum value. This constructs a consistent economy attaining each candidate feature. Compactness of $C_\eta$ gives its extrema. If its minimum fails the strict inequality, the constructed economy supplies a response violating strict superiority.
\end{proof}
The minimum identifies a guarantee, not a favorable selection from a tie. With interval messages the lifted formulation of Theorem~\ref{thm:oracle} retains the unknown query values and their supporting policies. Simply replacing all exact values in \eqref{eq:r19-exact-messages} by unrelated interval endpoints need not preserve sharpness.

\subsection{Known dynamics and the global response-loss budget}
When the finite dynamic model $M$ is known, let $X_M$ be its occupancy-measure polytope, $r^\top x$ private payoff, $Z^\top x$ response, and $W=\max_{x\in X_M}r^\top x$. Then the exact known-model response set is
\begin{equation}
 R_\eta(M)=\{Z^\top x:x\in X_M,\ r^\top x\geq W-\eta\}.
 \label{eq:r19-known-response}
\end{equation}
This is a different information set from \eqref{eq:r19-exact-messages}. The former uses transition feasibility; the latter allows every affine-response economy consistent with value messages. The performance-difference identity expresses the single restriction in \eqref{eq:r19-known-response} as a total expected Bellman-loss budget. It does not permit an $\eta$ loss at every state and date.

For a response direction $w$, finite-dimensional linear-programming duality gives
\begin{equation}
 h_{R_\eta(M)}(w)=\inf_{\lambda\geq0}\left\{\max_{x\in X_M}(w^\top Z^\top x+\lambda r^\top x)-\lambda(W-\eta)\right\}.
 \label{eq:r19-directed-dual}
\end{equation}
Writing $W(t)=\max_{x\in X_M}(r^\top x+t^\top Z^\top x)$, each $\lambda>0$ gives the bound $\lambda\{W(w/\lambda)-W(0)+\eta\}$. Certified values $U(w/\lambda)$ and $L(0)$ therefore give a response upper bound $\lambda\{U(w/\lambda)-L(0)+\eta\}$. Value uncertainty is amplified by $\lambda$ and must be charged before comparing purchasers. These occupancy and duality statements are established tools. The additional economic claim in Theorem~\ref{thm:r19-information} is the exact separation between an uninformative query family and a decision-identifying joint query, including its institutional and noise dependence.
